\documentclass[12pt,a4paper]{letter}
\usepackage{geometry}
\geometry{margin=2.5cm}
\usepackage{hyperref}

\signature{Christopher R. B. Wilson}
\address{[ADDRESS LINE 1]\\
         [ADDRESS LINE 2]\\
         [City, Country, Postcode]\\
         \texttt{C.Isomorph@gmail.com}\\
         ORCID: [ORCID-ID]}

\begin{document}

\begin{letter}{[Editor Name]\\
               Editor / Handling Editor\\
               Physical Review D\\
               American Physical Society}

\opening{Dear [Editor Name],}

I am pleased to submit for consideration by \textit{Physical Review D} the
manuscript:

\begin{center}
``Two Structural Obstructions to Galactic Dark Matter in Non-Minimally
Coupled Scalar-Tensor Gravity''
\end{center}

\textbf{Summary.}
The paper proves two no-go theorems constraining dark-matter production in
scalar-tensor extensions of general relativity with coupling $F(\Psi) =
\tfrac{1}{2} + \Lambda_0\Psi^2$.

\emph{Theorem~1 ($G_{\rm eff}$ ceiling).}
For any theory in this coupling class, the Solar-System Cassini
post-Newtonian bound ($|\gamma_{\rm PPN}-1| < 2.3\times10^{-5}$) limits
$\Lambda_0 < 1.29\times10^{-3}\,M_{\rm Pl}^{-2}$, giving a maximum
effective gravitational coupling $G_{\rm eff}/G_N \leq 1.018$.
Flat galactic rotation curves require $G_{\rm eff}/G_N \approx 3.1$, which
demands $\Lambda_0 \geq 0.154\,M_{\rm Pl}^{-2}$.  The two constraints are
incompatible by a factor of $\sim\!119$.

\emph{Theorem~2 (shape inversion).}
For any density-dependent coupling $\beta(\rho_b)$ that is monotonically
increasing in baryonic density, the predicted gravitational enhancement
$E_{\rm model}(r)$ tracks the baryon density and decreases with galactic
radius, while flat rotation curves require an enhancement that increases
with radius.  The anti-correlation is geometric and parameter-independent;
numerical evaluation on NGC~3198 gives Pearson $r = -0.81$.

Together the theorems close the two principal mechanism channels (effective
coupling modification and density-dependent fifth force) for dark matter
production in the quadratic non-minimal coupling class.  Curvature-memory
scalar-tensor gravity (CMSTG) is used as the concrete worked example.
The scope and limits of each theorem are characterised precisely.

\textbf{Significance.}
Both theorems are derived from general arguments (the Brans--Dicke
correspondence and monotonicity of density profiles) rather than
CMSTG-specific results, making them applicable to the broader family of
non-minimally coupled scalar-tensor theories.  Theorem~2 in particular
applies to any monotonically increasing chameleon or symmetron coupling
without further derivation.

\textbf{Relation to prior work.}
This paper is a companion short note to a four-paper CMSTG series submitted
to JCAP (Papers~I--IV; Paper~IV documents the worked examples for both
theorems).  The theorems are stated as general results here; the
JCAP paper uses them as steps in a broader enumeration of mechanism channels.

\textbf{Originality.}
All results are original.  The paper has not been submitted elsewhere.
Simulation code is publicly archived at
\url{https://github.com/cisomorph/Curvature-Memory_Scalar-Tensor_Gravity}.

\textbf{Suggested reviewers (optional):}\\
{[SUGGESTED REVIEWER 1: Name, institution, e-mail, expertise]}\\
{[SUGGESTED REVIEWER 2: Name, institution, e-mail, expertise]}\\
{[SUGGESTED REVIEWER 3: Name, institution, e-mail, expertise]}

I am available for correspondence at the address above.

\closing{Yours sincerely,}

\end{letter}
\end{document}
